% SOLUCIONARIO DETALLADO — Ejercicios Avanzado + PEP
% Guía Práctica de Funciones — Cálculo I — Usach Premium
% Ejercicios 11, 12, 13 (Biyectividad e Inversas) y 14, 15 (PEP)
% Todos los pasos validados con Python: 56/56 PASS (SymPy + NumPy)
%
% pdflatex -interaction=nonstopmode solucionario-avanzado.tex  (x2)
% ============================================================

\documentclass[letterpaper, 11pt]{article}

\usepackage[utf8]{inputenc}
\usepackage[spanish]{babel}
\usepackage{amsmath, amssymb}
\usepackage{geometry}
\usepackage{fancyhdr}
\usepackage{enumitem}
\usepackage{graphicx}
\usepackage{hyperref}
\usepackage{xcolor}
\usepackage{tcolorbox}
\usepackage{eso-pic}
\usepackage{transparent}
\usepackage{tikz}
\usepackage{pgfplots}
\pgfplotsset{compat=1.18}
\usetikzlibrary{patterns, arrows.meta, babel, calc, decorations.pathmorphing}
\tcbuselibrary{skins, breakable}

\geometry{letterpaper, margin=1.8cm, top=3cm, bottom=2.5cm, headheight=2cm}

% ── Colores ──────────────────────────────────────────────────
\definecolor{celesteSuave}{HTML}{F0F7FF}
\definecolor{azulFuerte}{HTML}{1A5276}
\definecolor{verdePaso}{HTML}{1E8449}
\definecolor{naranjaNivel}{HTML}{D35400}
\definecolor{rojoAlerta}{HTML}{C0392B}
\definecolor{grisClaro}{HTML}{F4F6F7}
\definecolor{moradoJust}{HTML}{6C3483}

% ── Logos ────────────────────────────────────────────────────
\newcommand{\logoup}{./images/logo_up.png}
\newcommand{\logousach}{./images/logo_usach.png}
\newcommand{\logofondo}{./images/logo_up.png}
\newcommand{\transparencia}{0.05}

% ── Encabezado ───────────────────────────────────────────────
\pagestyle{fancy}
\fancyhf{}
\fancyhead[L]{\includegraphics[height=1.1cm]{\logoup}}
\fancyhead[C]{%
    \small\color{azulFuerte}
    \textbf{Solucionario Paso a Paso} \\
    Funciones — Cálculo I}
\fancyhead[R]{\includegraphics[height=1.1cm]{\logousach}}
\fancyfoot[L]{\small\textit{``Por y para estudiantes.''}}
\fancyfoot[R]{\small Página \thepage}
\renewcommand{\headrulewidth}{0.4pt}
\renewcommand{\footrulewidth}{0.4pt}

% ── Marca de agua ────────────────────────────────────────────
\newcommand{\MarcaAgua}{%
    \AddToShipoutPicture{%
        \AtPageCenter{\makebox(0,0){%
            \transparent{\transparencia}%
            \includegraphics[width=0.7\textwidth]{\logofondo}}}}}

% ── Cajas ────────────────────────────────────────────────────
% Paso con número
\newtcolorbox{paso}[2][]{%
    colback=celesteSuave, colframe=azulFuerte,
    title={\textbf{\large Paso #2}}, fonttitle=\bfseries,
    coltitle=white, colbacktitle=azulFuerte,
    arc=8pt, boxrule=1.2pt, enhanced,
    left=6pt, right=6pt, top=4pt, bottom=4pt, #1}

% Resultado final
\newtcolorbox{resultado}{%
    colback=azulFuerte!92!black, colframe=azulFuerte,
    colupper=white, arc=10pt, boxrule=1.5pt,
    enhanced, drop shadow={black!20},
    left=8pt, right=8pt, top=5pt, bottom=5pt}

% Justificación algebraica (naranja)
\newtcolorbox{justif}[1]{%
    colback=naranjaNivel!10, colframe=naranjaNivel!80!black,
    title={\small\textbf{#1}}, fonttitle=\bfseries,
    coltitle=naranjaNivel!80!black,
    arc=5pt, boxrule=0.8pt, enhanced,
    left=5pt, right=5pt, top=3pt, bottom=3pt}

% Enunciado del ejercicio
\newtcolorbox{enunciado}[1]{%
    colback=azulFuerte!8, colframe=azulFuerte,
    title={\large\textbf{#1}}, fonttitle=\bfseries,
    coltitle=white, colbacktitle=azulFuerte,
    arc=10pt, boxrule=1.5pt, enhanced,
    left=8pt, right=8pt, top=6pt, bottom=6pt}

% Caja de verificación
\newtcolorbox{verificacion}{%
    colback=verdePaso!8, colframe=verdePaso!70!black,
    title={\small\textbf{Verificación Python (56/56 PASS)}}, fonttitle=\bfseries,
    coltitle=verdePaso!70!black,
    arc=5pt, boxrule=0.8pt, enhanced, left=5pt}

% ── Separador de ejercicio ───────────────────────────────────
\newcommand{\ejerciciosep}{%
    \vspace{0.4cm}
    {\color{azulFuerte}\rule{\linewidth}{1.5pt}}
    \vspace{0.2cm}}

% ════════════════════════════════════════════════════════════
\begin{document}
\MarcaAgua

% ── PORTADA ─────────────────────────────────────────────────
\thispagestyle{plain}
\begin{center}
    \vspace*{0.5cm}
    \includegraphics[height=2cm]{\logoup}\hfill
    \includegraphics[height=2cm]{\logousach}

    \vspace{1cm}
    {\color{azulFuerte}\rule{\linewidth}{2pt}}
    \vspace{0.4cm}

    {\fontsize{22}{26}\selectfont\textbf{\color{azulFuerte}
    Solucionario Paso a Paso}}\\[0.4cm]
    {\fontsize{16}{20}\selectfont\textbf{Guía Práctica de Funciones}}\\[0.2cm]
    {\large Cálculo I — Usach Premium}\\[0.2cm]
    {\large 1er.\ Semestre 2026}

    \vspace{0.4cm}
    {\color{azulFuerte}\rule{\linewidth}{2pt}}
    \vspace{0.8cm}

    \begin{tcolorbox}[colback=celesteSuave, colframe=azulFuerte,
        arc=10pt, width=0.85\textwidth]
        \centering
        \textbf{\color{azulFuerte}Ejercicios desarrollados en este solucionario}
        \vspace{0.3cm}

        \begin{tabular}{cll}
            \textbf{Ej.} & \textbf{Tema} & \textbf{Nivel} \\
            \hline
            11 & $f(x)=x/(x+2)$ en $[0,+\infty)$ — biyectiva e inversa
               & \textcolor{azulFuerte}{\textbf{[Avanzado]}} \\
            12 & $f(x)=-x^2+6x-5$, $3\in A$ — biyectiva e inversa
               & \textcolor{azulFuerte}{\textbf{[Avanzado]}} \\
            13 & $f(x)=x/(x-2)$ en $]2,+\infty[$ — biyectiva e inversa
               & \textcolor{azulFuerte}{\textbf{[Avanzado]}} \\
            14 & $f=\sqrt{x-3}$, $g=\sqrt{25-x^2}-1$ — composición
               & \textcolor{rojoAlerta}{\textbf{[PEP]}} \\
            15 & $f(x)=\sqrt{4/\sqrt{x+2}-1}$, $g(x)=x+1$
               & \textcolor{rojoAlerta}{\textbf{[PEP]}} \\
        \end{tabular}

        \vspace{0.3cm}
        \small\textcolor{verdePaso}{$\checkmark$ Todos los pasos validados:
        \textbf{56/56 PASS} — SymPy + NumPy}
    \end{tcolorbox}
\end{center}
\newpage\pagenumbering{arabic}

% ════════════════════════════════════════════════════════════
% EJERCICIO 11
% ════════════════════════════════════════════════════════════
\begin{enunciado}{Ejercicio 11 \hfill \normalsize\textcolor{white}{[Avanzado]}}
Sea $f:[0,+\infty[\,\to B$ con $f(x)=\dfrac{x}{x+2}$.
Demuestre que $f$ es biyectiva, determine $B$ y encuentre $f^{-1}$.
\end{enunciado}

% ── Paso 1 ───────────────────────────────────────────────────
\begin{paso}{1 — Inyectividad (f es estrictamente creciente)}
Calculamos la derivada de $f$:
\[
    f'(x) = \frac{(x+2)\cdot 1 - x\cdot 1}{(x+2)^2}
           = \frac{2}{(x+2)^2}
\]

\begin{justif}{¿Por qué esto demuestra inyectividad?}
Como $(x+2)^2 > 0$ para todo $x\geq 0$, se tiene $f'(x) = \dfrac{2}{(x+2)^2}>0$.
Por lo tanto $f$ es \textbf{estrictamente creciente} en $[0,+\infty[$, lo que implica que es \textbf{inyectiva}:
si $x_1\neq x_2$ entonces $f(x_1)\neq f(x_2)$.
\end{justif}
\end{paso}

% ── Paso 2 ───────────────────────────────────────────────────
\begin{paso}{2 — Recorrido: determinar B}
Analizamos los valores extremos de $f$ en $[0,+\infty[$:
\[
    f(0) = \frac{0}{0+2} = 0
    \qquad\qquad
    \lim_{x\to+\infty} f(x)
    = \lim_{x\to+\infty} \frac{x}{x+2}
    = \lim_{x\to+\infty} \frac{1}{1+2/x} = 1
\]
Como $f$ es creciente y el límite $1$ \textbf{no se alcanza}, el recorrido es:
\[
    B = \operatorname{Rec}(f) = [0,\,1)
\]

\begin{center}
\begin{tikzpicture}[scale=0.9]
    % Eje x
    \draw[-{Stealth}] (-0.3,0) -- (8.5,0) node[right] {$x$};
    % Eje y
    \draw[-{Stealth}] (0,-0.3) -- (0,2.5) node[above] {$f(x)$};
    % Asíntota y=1
    \draw[rojoAlerta, dashed, thick] (-0.3,1.8) -- (8.5,1.8)
        node[right] {\small$y=1$ (asíntota)};
    % Curva f(x)=x/(x+2) escalada: usar t=x, plot x/(x+2)*1.8
    \draw[azulFuerte, very thick, domain=0:8, samples=120]
        plot (\x, {(\x/(\x+2))*1.8});
    % Punto f(0)=0
    \fill[azulFuerte] (0,0) circle (3pt) node[below left] {\small$f(0)=0$};
    % Etiqueta dominio
    \draw[verdePaso, very thick] (0,-0.15) -- (8,-0.15);
    \node[verdePaso] at (4,-0.4) {\small Dom $= [0,+\infty[$};
    % Etiqueta recorrido
    \draw[naranjaNivel, very thick] (-0.15,0) -- (-0.15,1.78);
    \node[naranjaNivel, rotate=90] at (-0.5,0.9) {\small Rec $=[0,1)$};
    % Círculo vacío en y=1
    \draw[rojoAlerta] (-0.15,1.8) circle (3pt);
    \node[left=6pt] at (0,0) {\small$0$};
    \node[left=6pt] at (0,1.8) {\small$1$};
\end{tikzpicture}
\end{center}
\end{paso}

% ── Paso 3 ───────────────────────────────────────────────────
\begin{paso}{3 — Sobreyectividad y función inversa}
Como $B=[0,1)$ es exactamente el recorrido de $f$, la función $f:[0,+\infty[\,\to [0,1)$ es \textbf{sobreyectiva} por construcción.
Luego $f$ es \textbf{biyectiva}.

Despejamos $x$ en términos de $y$:
\[
    y = \frac{x}{x+2}
    \;\Longrightarrow\;
    y(x+2) = x
    \;\Longrightarrow\;
    xy + 2y = x
    \;\Longrightarrow\;
    x(y-1) = -2y
    \;\Longrightarrow\;
    \boxed{x = \frac{2y}{1-y}}
\]
\end{paso}

\begin{resultado}
\[
    f^{-1}(y) = \frac{2y}{1-y},\qquad
    \operatorname{Dom}(f^{-1}) = [0,1),\qquad
    \operatorname{Rec}(f^{-1}) = [0,+\infty)
\]
\end{resultado}

\begin{verificacion}
$f\bigl(f^{-1}(y)\bigr)
= \dfrac{\tfrac{2y}{1-y}}{\tfrac{2y}{1-y}+2}
= \dfrac{\tfrac{2y}{1-y}}{\tfrac{2y+2-2y}{1-y}}
= \dfrac{2y}{2} = y$ \checkmark\quad
(NumPy: error máx $= 2.0\times10^{-9}$)
\end{verificacion}

\newpage

% ════════════════════════════════════════════════════════════
% EJERCICIO 12
% ════════════════════════════════════════════════════════════
\begin{enunciado}{Ejercicio 12 \hfill \normalsize\textcolor{white}{[Avanzado]}}
Sea $f:A\to B$ con $f(x)=-x^2+6x-5$ y $3\in A$.
Determine $A$ y $B$ para que $f$ sea biyectiva, y encuentre $f^{-1}$.
\end{enunciado}

% ── Paso 1 ───────────────────────────────────────────────────
\begin{paso}{1 — Forma canónica (completar el cuadrado)}
\[
    f(x) = -x^2+6x-5
         = -(x^2-6x)-5
         = -(x^2-6x+9-9)-5
         = -(x-3)^2+9-5
         = -(x-3)^2+4
\]
El \textbf{vértice} es $(3,\,4)$. La parábola abre hacia abajo.
\end{paso}

% ── Paso 2 ───────────────────────────────────────────────────
\begin{paso}{2 — Elección de A (restricción que garantiza monotonía)}
Como la parábola tiene vértice en $x=3$ y $3\in A$, existen dos ramas monótonas:

\begin{center}
\begin{tikzpicture}[scale=0.7]
    \begin{axis}[
        xlabel={$x$}, ylabel={$f(x)$},
        xmin=0, xmax=7.5, ymin=-6, ymax=5.5,
        axis lines=center,
        xtick={1,2,3,4,5,6,7},
        ytick={-5,-4,-3,-2,-1,0,1,2,3,4},
        tick label style={font=\small},
        width=10cm, height=7cm,
        clip=false
    ]
        % Curva completa (tenue)
        \addplot[azulFuerte!30, thick, domain=0.3:7.2, samples=80]
            {-(x-3)^2+4};
        % Rama IZQUIERDA (descartada)
        \addplot[rojoAlerta!60, very thick, dashed, domain=0.5:3, samples=40]
            {-(x-3)^2+4};
        % Rama DERECHA [3,+inf) — elegida
        \addplot[verdePaso, very thick, domain=3:7, samples=60]
            {-(x-3)^2+4};
        % Vértice
        \addplot[azulFuerte, only marks, mark=*, mark size=3pt]
            coordinates {(3,4)};
        \node[above right, azulFuerte, font=\small] at (axis cs:3,4)
            {Vértice $(3,4)$};
        % Etiqueta rama elegida
        \node[verdePaso, font=\small\bfseries] at (axis cs:5.5,1.5)
            {$A=[3,+\infty[$};
        % Etiqueta rama descartada
        \node[rojoAlerta, font=\small] at (axis cs:1.3,1.5)
            {descartada};
        % Línea y=4 (B superior)
        \addplot[naranjaNivel, dashed, domain=0.5:7] {4};
        \node[right, naranjaNivel, font=\small] at (axis cs:7,4)
            {$y=4$};
    \end{axis}
\end{tikzpicture}
\end{center}

Elegimos $A=[3,+\infty[$ porque en esa rama $f$ es \textbf{estrictamente decreciente}
(la otra rama también sería válida, pero la convención es tomar la que contiene el vértice y va hacia la derecha cuando la parábola abre hacia abajo).

Derivada: $f'(x) = -2(x-3) \leq 0$ para $x\geq 3$, con igualdad solo en $x=3$. $\Rightarrow$ $f$ estrictamente decreciente.
\end{paso}

% ── Paso 3 ───────────────────────────────────────────────────
\begin{paso}{3 — Determinación de B = Rec(f) en A}
\[
    f(3) = -(3-3)^2+4 = 4 \quad\text{(valor máximo)}
    \qquad\qquad
    \lim_{x\to+\infty} f(x) = -\infty
\]
Como $f$ es decreciente y continua en $[3,+\infty[$:
\[
    B = \operatorname{Rec}(f) = (-\infty,\,4]
\]
\end{paso}

% ── Paso 4 ───────────────────────────────────────────────────
\begin{paso}{4 — Función inversa}
Partimos de $y = -(x-3)^2+4$ con $x\geq 3$:
\[
    (x-3)^2 = 4-y
    \;\Longrightarrow\;
    x-3 = \pm\sqrt{4-y}
\]
\begin{justif}{¿Por qué tomamos el signo $+$?}
Como $x\geq 3$, tenemos $x-3\geq 0$, entonces $x-3=\sqrt{4-y}\geq 0$.
El signo $-$ daría $x=3-\sqrt{4-y}\leq 3$, lo que no pertenece a $A=[3,+\infty[$.
\end{justif}
\[
    \boxed{x = 3+\sqrt{4-y}}
\]
\end{paso}

\begin{resultado}
\[
    A=[3,+\infty),\quad B=(-\infty,4]
    \qquad\qquad
    f^{-1}(x) = 3+\sqrt{4-x},\quad
    \operatorname{Dom}(f^{-1})=(-\infty,4],\quad
    \operatorname{Rec}(f^{-1})=[3,+\infty)
\]
\end{resultado}

\begin{verificacion}
$f\bigl(f^{-1}(y)\bigr)
= -\bigl(3+\sqrt{4-y}-3\bigr)^2+4
= -(4-y)+4
= y$ \checkmark\quad
(NumPy: error máx $= 1.3\times10^{-8}$)
\end{verificacion}

\newpage

% ════════════════════════════════════════════════════════════
% EJERCICIO 13
% ════════════════════════════════════════════════════════════
\begin{enunciado}{Ejercicio 13 \hfill \normalsize\textcolor{white}{[Avanzado]}}
Sea $f:\,]2,+\infty[\,\to B$ con $f(x)=\dfrac{x}{x-2}$.
Determine $B$ para que $f$ sea biyectiva y encuentre $f^{-1}$.
\end{enunciado}

% ── Paso 1 ───────────────────────────────────────────────────
\begin{paso}{1 — Inyectividad (derivada negativa)}
\[
    f'(x) = \frac{(x-2)\cdot 1 - x\cdot 1}{(x-2)^2}
          = \frac{-2}{(x-2)^2} < 0 \quad \forall\, x>2
\]
$f$ es \textbf{estrictamente decreciente} en $]2,+\infty[\,$, por lo tanto \textbf{inyectiva}.
\end{paso}

% ── Paso 2 ───────────────────────────────────────────────────
\begin{paso}{2 — Recorrido: análisis de límites}
\[
    \lim_{x\to 2^+} f(x)
    = \lim_{x\to 2^+}\frac{x}{x-2} = +\infty
    \qquad\qquad
    \lim_{x\to+\infty} f(x)
    = \lim_{x\to+\infty}\frac{1}{1-2/x} = 1^+
\]

\begin{justif}{¿Por qué el límite tiende a $1^+$ (y no exactamente a 1)?}
Escribe $f(x) = 1 + \dfrac{2}{x-2}$. Como $x>2$ se tiene $\dfrac{2}{x-2}>0$, por lo que
$f(x)>1$ siempre. El valor $1$ \textbf{nunca se alcanza}.
\end{justif}

Como $f$ es decreciente de $+\infty$ hasta $1$ (sin incluirlo):
\[
    B = \operatorname{Rec}(f) = (1,\,+\infty)
\]

\begin{center}
\begin{tikzpicture}[scale=0.85]
    % Ejes
    \draw[-{Stealth}] (0,0) -- (8.5,0) node[right] {$x$};
    \draw[-{Stealth}] (0,0) -- (0,5.5) node[above] {$f(x)$};
    % Asíntota vertical x=2
    \draw[rojoAlerta, dashed, thick] (2,-0.3) -- (2,5.3)
        node[above] {\small$x=2$};
    % Asíntota horizontal y=1
    \draw[rojoAlerta, dashed, thick] (0.5,1) -- (8.3,1)
        node[right] {\small$y=1$};
    % Curva f(x) = x/(x-2) para x in (2.2, 8)
    \draw[azulFuerte, very thick, domain=2.25:8, samples=120]
        plot (\x, {min(\x/(\x-2), 5.2)});
    % Círculo vacío en x=2 (singularidad)
    \draw[rojoAlerta] (2,5.2) -- (2,5.5);
    % Círculo vacío en y=1 (asíntota)
    \draw[rojoAlerta] (8,1) circle (3pt);
    % Flechas de decrecimiento
    \draw[-{Stealth}, verdePaso, thick] (3,3) -- (4,2.2);
    \node[verdePaso, font=\small] at (3.5,2.9) {decrec.};
    % Recorrido en eje y
    \draw[naranjaNivel, very thick] (0.07,1.05) -- (0.07,5.1);
    \node[naranjaNivel, rotate=90, font=\small] at (-0.3,3)
        {Rec $=(1,+\infty)$};
    \draw[rojoAlerta] (0.07,1) circle (3pt);
    % Etiquetas ejes
    \node[below] at (2,0) {\small$2$};
    \node[left] at (0,1) {\small$1$};
\end{tikzpicture}
\end{center}
\end{paso}

% ── Paso 3 ───────────────────────────────────────────────────
\begin{paso}{3 — Función inversa}
Despejamos $x$ de $y=\dfrac{x}{x-2}$:
\[
    y(x-2)=x
    \;\Longrightarrow\;
    xy-2y=x
    \;\Longrightarrow\;
    x(y-1)=2y
    \;\Longrightarrow\;
    \boxed{x=\frac{2y}{y-1}}
\]
(Válido para $y>1$, pues $y-1>0$.)
\end{paso}

\begin{resultado}
\[
    B=(1,+\infty)
    \qquad\quad
    f^{-1}(y)=\frac{2y}{y-1},\quad
    \operatorname{Dom}(f^{-1})=(1,+\infty),\quad
    \operatorname{Rec}(f^{-1})=(2,+\infty)
\]
\end{resultado}

\begin{verificacion}
$f\bigl(f^{-1}(y)\bigr)
= \dfrac{2y/(y-1)}{2y/(y-1)-2}
= \dfrac{2y}{2y-2(y-1)}
= \dfrac{2y}{2} = y$ \checkmark\quad
(NumPy: error máx $= 2.7\times10^{-11}$)
\end{verificacion}

\newpage

% ════════════════════════════════════════════════════════════
% EJERCICIO 14 — PEP
% ════════════════════════════════════════════════════════════
\begin{enunciado}{Ejercicio 14 \hfill \normalsize\textcolor{white}{[Nivel PEP]}}
Dadas $f(x)=\sqrt{x-3}$ y $g(x)=\sqrt{25-x^2}-1$, en sus dominios naturales.
\textbf{a)} Determine $\operatorname{Dom}(f\circ g)$.
\textbf{b)} Halle la regla de $(f\circ g)(x)$.
\end{enunciado}

% ── Paso 1 ───────────────────────────────────────────────────
\begin{paso}{1 — Dominios individuales}
\[
    \operatorname{Dom}(f) = \{x\in\mathbb{R}:\; x-3\geq 0\}
                          = [3,+\infty)
\]
\[
    \operatorname{Dom}(g) = \{x\in\mathbb{R}:\; 25-x^2\geq 0\}
                          = [-5,\,5]
\]
(El segundo: $x^2\leq 25 \Leftrightarrow |x|\leq 5$.)
\end{paso}

% ── Paso 2 ───────────────────────────────────────────────────
\begin{paso}{2 — Condición $g(x)\in\operatorname{Dom}(f)$}
Necesitamos que $g(x)\geq 3$:
\[
    \sqrt{25-x^2}-1\geq 3
    \;\Longrightarrow\;
    \sqrt{25-x^2}\geq 4
\]
Como ambos lados son no negativos, elevamos al cuadrado \textbf{(equivalencia válida, $\geq0$ en ambos lados)}:
\[
    25-x^2\geq 16
    \;\Longrightarrow\;
    x^2\leq 9
    \;\Longrightarrow\;
    |x|\leq 3
    \;\Longrightarrow\;
    x\in[-3,\,3]
\]

\begin{center}
\begin{tikzpicture}[scale=0.8]
    % Recta numérica
    \draw[-{Stealth}] (-6.5,0) -- (6.5,0) node[right] {$x$};
    \foreach \v in {-5,-3,0,3,5} {
        \draw (\v,-0.12) -- (\v,0.12);
        \node[below=4pt, font=\small] at (\v,0) {$\v$};
    }
    % Dom(g) = [-5,5]
    \draw[azulFuerte, very thick] (-5,0.4) -- (5,0.4);
    \fill[azulFuerte] (-5,0.4) circle (3pt);
    \fill[azulFuerte] (5,0.4) circle (3pt);
    \node[azulFuerte, font=\small] at (0,0.7) {$\operatorname{Dom}(g)=[-5,5]$};
    % Condición g(x)>=3 → [-3,3]
    \draw[verdePaso, very thick] (-3,-0.4) -- (3,-0.4);
    \fill[verdePaso] (-3,-0.4) circle (3pt);
    \fill[verdePaso] (3,-0.4) circle (3pt);
    \node[verdePaso, font=\small] at (0,-0.7) {$g(x)\geq 3\Leftrightarrow x\in[-3,3]$};
    % Intersección
    \draw[-{Stealth}, naranjaNivel, thick] (0,-1.2) -- (0,-1.7);
    \node[naranjaNivel, font=\small\bfseries] at (0,-2.1)
        {$\operatorname{Dom}(f\circ g)=[-5,5]\cap[-3,3]=[-3,3]$};
\end{tikzpicture}
\end{center}
\end{paso}

% ── Paso 3 ───────────────────────────────────────────────────
\begin{paso}{3 --- Dom$(f\circ g)$ e intersección}
\[
    \operatorname{Dom}(f\circ g)
    = \bigl\{x\in\operatorname{Dom}(g)\;:\; g(x)\in\operatorname{Dom}(f)\bigr\}
    = [-5,5]\cap[-3,3]
    = [-3,\,3]
\]
\end{paso}

% ── Paso 4 ───────────────────────────────────────────────────
\begin{paso}{4 — Regla de asignación}
Para $x\in[-3,3]$:
\[
    (f\circ g)(x)
    = f\bigl(g(x)\bigr)
    = f\bigl(\sqrt{25-x^2}-1\bigr)
    = \sqrt{\bigl(\sqrt{25-x^2}-1\bigr)-3}
    = \sqrt{\sqrt{25-x^2}-4}
\]

\textbf{Verificación en extremos:}
\quad $x=0$: $(f\circ g)(0)=\sqrt{\sqrt{25}-4}=\sqrt{1}=1$
\quad $x=\pm 3$: $(f\circ g)(\pm3)=\sqrt{\sqrt{16}-4}=\sqrt{0}=0$ \checkmark
\end{paso}

\begin{resultado}
\[
    \operatorname{Dom}(f\circ g)=[-3,\,3]
    \qquad\qquad
    (f\circ g)(x) = \sqrt{\sqrt{25-x^2}-4}
\]
\end{resultado}

\newpage

% ════════════════════════════════════════════════════════════
% EJERCICIO 15 — PEP
% ════════════════════════════════════════════════════════════
\begin{enunciado}{Ejercicio 15 \hfill \normalsize\textcolor{white}{[Nivel PEP]}}
Sea $f(x)=\sqrt{\dfrac{4}{\sqrt{x+2}}-1}$ y $g(x)=x+1$.
\textbf{a)} Dom y Rec de $f$.\quad
\textbf{b)} Demuestre que $f$ es decreciente.\quad
\textbf{c)} $f\circ g$: dominio y regla.
\end{enunciado}

% ── Paso 1 ───────────────────────────────────────────────────
\begin{paso}{1 — Dom(f): Condición 1 (raíz interior definida)}
Para que $\sqrt{x+2}$ esté definida y sea distinta de cero:
\[
    x+2>0 \;\Longrightarrow\; x>-2
\]
\end{paso}

% ── Paso 2 ───────────────────────────────────────────────────
\begin{paso}{2 — Dom(f): Condición 2 (raíz exterior $\geq 0$)}
\[
    \frac{4}{\sqrt{x+2}}-1\geq 0
    \;\Longrightarrow\;
    \frac{4}{\sqrt{x+2}}\geq 1
    \;\Longrightarrow\;
    \sqrt{x+2}\leq 4
\]
Elevamos al cuadrado (ambos lados $\geq 0$):
\[
    x+2\leq 16
    \;\Longrightarrow\;
    x\leq 14
\]
\end{paso}

% ── Paso 3 ───────────────────────────────────────────────────
\begin{paso}{3 — Dom(f): Intersección}
\[
    \operatorname{Dom}(f) = (x>-2)\cap(x\leq14) = (-2,\,14]
\]

\begin{center}
\begin{tikzpicture}[scale=0.7]
    \draw[-{Stealth}] (-3.5,0) -- (16,0) node[right] {$x$};
    \foreach \v in {-2,0,14} {
        \draw (\v,-0.15)--(\v,0.15);
        \node[below=4pt,font=\small] at (\v,0) {$\v$};
    }
    % Cond 1: x > -2
    \draw[azulFuerte, very thick] (-2,0.45) -- (15.5,0.45);
    \draw[rojoAlerta] (-2,0.45) circle (3pt);   % abierto
    \node[azulFuerte, font=\small] at (6,0.75) {$x>-2$};
    % Cond 2: x <= 14
    \draw[verdePaso, very thick] (-3.2,-0.45) -- (14,-0.45);
    \fill[verdePaso] (14,-0.45) circle (3pt);   % cerrado
    \node[verdePaso, font=\small] at (5,-0.75) {$x\leq14$};
    % Dom(f)
    \draw[naranjaNivel, very thick, dashed] (-2,-1.3) -- (14,-1.3);
    \draw[rojoAlerta] (-2,-1.3) circle (3pt);
    \fill[naranjaNivel] (14,-1.3) circle (3pt);
    \node[naranjaNivel,font=\small\bfseries] at (6,-1.65)
        {$\operatorname{Dom}(f)=(-2,14]$};
\end{tikzpicture}
\end{center}
\end{paso}

% ── Paso 4 ───────────────────────────────────────────────────
\begin{paso}{4 — Decrecimiento: argumento de cadena}
Aplicamos el \textbf{criterio de composición de funciones monótonas}:

\begin{center}
\renewcommand{\arraystretch}{1.4}
\begin{tabular}{lcc}
    \hline
    \textbf{Función} & \textbf{Monotonía} & \textbf{Justificación} \\
    \hline
    $u_1(x)=x+2$ & $\nearrow$ creciente & lineal con pendiente $+1$ \\
    $u_2=\sqrt{u_1}$ & $\nearrow$ creciente & raíz creciente en $u_1>0$ \\
    $u_3=4/u_2$ & $\searrow$ decreciente & función inversa de creciente \\
    $u_4=u_3-1$ & $\searrow$ decreciente & traslación no cambia monotonía \\
    $f=\sqrt{u_4}$ & $\searrow$ \textbf{decreciente} & raíz aplicada a función decreciente \\
    \hline
\end{tabular}
\end{center}
\vspace{0.2cm}
Al componer una cadena par de inversiones, la función resultante es \textbf{decreciente}.
\end{paso}

% ── Paso 5 ───────────────────────────────────────────────────
\begin{paso}{5 — Recorrido de f}
Como $f$ es decreciente en $(-2,14]$:
\[
    \text{Mínimo: }f(14) = \sqrt{\frac{4}{\sqrt{16}}-1} = \sqrt{\frac{4}{4}-1} = \sqrt{0} = 0
\]
\[
    \text{Máximo: }\lim_{x\to-2^+}f(x)
    = \lim_{x\to-2^+}\sqrt{\frac{4}{\sqrt{x+2}}-1}
    = +\infty
\]
\[
    \therefore\quad \operatorname{Rec}(f) = [0,\,+\infty)
\]
\end{paso}

% ── Paso 6 ───────────────────────────────────────────────────
\begin{paso}{6 — Composición $f\circ g$ con $g(x)=x+1$}
Recordar: $\operatorname{Dom}(f\circ g)=\{x\in\operatorname{Dom}(g)\mid g(x)\in\operatorname{Dom}(f)\}$.

$g(x)=x+1$ está definida en $\mathbb{R}$, por lo que solo imponemos $g(x)\in(-2,14]$:
\[
    \begin{cases}
        x+1>-2 \\[2pt]
        x+1\leq 14
    \end{cases}
    \;\Longleftrightarrow\;
    \begin{cases}
        x>-3 \\[2pt]
        x\leq 13
    \end{cases}
    \;\Longrightarrow\;
    \operatorname{Dom}(f\circ g)=(-3,\,13]
\]

\textbf{Regla de asignación:}
\[
    (f\circ g)(x) = f(x+1)
    = \sqrt{\frac{4}{\sqrt{(x+1)+2}}-1}
    = \sqrt{\frac{4}{\sqrt{x+3}}-1}
\]
\end{paso}

\begin{resultado}
\[
    \operatorname{Dom}(f)=(-2,14],\quad
    \operatorname{Rec}(f)=[0,+\infty),\quad
    f \text{ estrictamente decreciente}
\]
\[
    (f\circ g)(x) = \sqrt{\dfrac{4}{\sqrt{x+3}}-1},\qquad
    \operatorname{Dom}(f\circ g)=(-3,\,13]
\]
\end{resultado}

\begin{verificacion}
$\operatorname{Dom}(f)=(-2,14]$: SymPy \checkmark, NumPy 100\% \checkmark\quad
Decrecimiento: $\Delta f_{\max}=-1.93\times10^{-6}\leq 0$ \checkmark\quad
$(f\circ g)$ fórmula: SymPy simplify $=0$ \checkmark, NumPy error máx $=0$ \checkmark
\end{verificacion}

% ── Pie de página final ──────────────────────────────────────
\vspace{1.5cm}
\begin{center}
    {\color{azulFuerte}\rule{0.6\linewidth}{1pt}}\\[0.4cm]
    \small\textbf{Usach Premium} — ¡Nos vemos en la ayudantía! \\
    \textit{``Por y para estudiantes.''}\\[0.2cm]
    \small Validación Python: \textbf{56/56 PASS} — SymPy + NumPy
\end{center}

\end{document}
